\section{Protocol Robustness and Integrity}
\label{sec:robust}

\begin{guidance}[Protocol robustness and integrity (10 marks)]
Show \textbf{framing}, \textbf{specific NACKs for malformed input}, and
\textbf{graceful recovery} from partial or garbage data. A table that maps each
kind of bad input to your mechanism, its response and the test that checks it
works well. Also explain why the server can never block forever (bounded
loops, timeouts, no blocking waits in the main loop). Give measured numbers
where you have them, e.g. the timeout you observed.
\end{guidance}

\ph{One-sentence summary of how your server handles bad input.}

\begin{table}[t]
\centering
\caption{\ph{Bad input, the server's handling, and the test that checks it.}}
\label{tab:threats}
\scriptsize
\begin{tabular}{@{}L{1.7cm}L{3.1cm}L{3.2cm}@{}}
\toprule
Bad input & Handling and response & Test \\
\midrule
Partial frame & \ph{\ldots} & \ph{\ldots} \\
Oversized frame & \ph{\ldots} & \ph{\ldots} \\
Garbage bytes & \ph{\ldots} & \ph{\ldots} \\
Wrong case or unknown command & \ph{\ldots} & \ph{\ldots} \\
Corrupted bytes (checksum) & \ph{\ldots} & \ph{\ldots} \\
UART errors (overrun, framing) & \ph{\ldots} & \ph{\ldots} \\
\bottomrule
\end{tabular}
\end{table}

\subsection{\ph{Heading about partial frames and timeouts}}
\ph{Your timeout rule, why you chose that value, and what you measured.}

\subsection{\ph{Heading about garbage and oversized input}}
\ph{How the server discards bad bytes and resynchronises on the next valid
frame.}

\subsection{\ph{Heading about why the server never hangs}}
\ph{Bounded loops, bounded transmit times, no blocking waits, watchdog (if
used).}
